%==============================================================================
% Section 8: Conclusion
%==============================================================================
\section{Conclusion}
\label{sec:conclusion}

In this paper we introduced the Weyl-squared coupling $\alpha C^2_\EC$ based on
the Einstein--Cartan connection into a homogeneous minisuperspace, compared it
across the three topologies $\Sthree\times\Sone$, $\Tthree\times\Sone$, and
$\Nilthree\times\Sone$, and organised what the EC--Weyl extension preserves and
what it updates in the mode dictionary of the homogeneous bulk geometric sector.

First, we confirmed that the new contribution of the EC--Weyl coupling does not
appear for an arbitrary torsion background; it drops out in the AX/VT
backgrounds and only essential modifications appear in the MX background.
Simultaneously, the Palatini-type non-propagation is preserved after the EC
extension, and the torsion sector does not develop as a kinetic mode.  This
shows that the basic structure of the homogeneous sector established in the LC
version does not break down substantially under the EC connection.

Second, for $\Sthree\times\Sone$, taking the homogeneous mode dictionary
established in paper~III as the basis, we made explicit which parts are
invariant and which receive corrections under the EC--Weyl extension.  In the
viewpoint of this paper, $\Sthree$ is still the reference topology providing the
richest homogeneous bulk geometry, and the EC extension is positioned not as
destroying its dictionary but as giving update rules.

Third, for $\Tthree\times\Sone$, the mode dictionary closes as essentially a
null/trivial structure.  That is, even after performing the EC--Weyl extension,
$\Tthree$ provides no new non-trivial mass structure in the homogeneous bulk
sector, and part of the cross-terms is also understood as having a purely
kinematic origin.  In this sense, $\Tthree$ functions as the reference
background that makes explicit ``what does not happen.''

Fourth, for $\Nilthree\times\Sone$, we showed that the EC--Weyl extension
produces a topology-specific EC slice minimum on the $\eta=V=0$ section and an
anisotropic split structure, and that the spin-2 quintet splitting and the
spin-1 uniaxial mass splitting are geometrically fixed.  Therefore $\Nilthree$
is positioned as the background carrying the EC-specific geometric phase
structure that is invisible in the LC version.

Taken together, this paper \textbf{closes the homogeneous three-topology bulk
sector under EC--Weyl, within the axial + vector-trace torsion ansatz}.  What
has been closed is the comparable homogeneous geometric sector based on the
common ansatz, and within this framework the division of roles among $\Sthree$,
$\Tthree$, and $\Nilthree$ has been clarified.

The results of this paper provide, within the adopted homogeneous torsion
ansatz, the foundation for subsequent theoretical development toward the
inhomogeneous sector and boundary responses.
